\begin{definition}
Let $g\in\mathsf{IncSeq}$ and let $R$ be a binary relation on a discrete
type $A$.  We call $(A,R)$ a $g$-better-relation when the following two
equivalent conditions hold:
\begin{enumerate}
\item for every locally constant $\varphi:\mathsf{IncSeq}\to A$ there are
an $f\in\mathsf{IncSeq}$ and a continuous homomorphism
\[
H:(\mathsf{IncSeq},\mathsf{RightComp}(g))\to(A,R)
\]
such that $H(x)=\varphi(f\circ x)$ for every $x$;
\item there is no continuous homomorphism
\[
(\mathsf{IncSeq},\mathsf{RightComp}(g))
\longrightarrow(A,R^{\complement}),
\]
where $R^{\complement}(a,b)$ means $\neg R(a,b)$.
\end{enumerate}
When $R$ is a quasi-order on $Q$, this is the assertion that $Q$ is
$g$-better-quasi-ordered.
\end{definition}
